%% sec-8-discussion.tex -- Discussion & Future Directions
\section{Discussion and Future Directions}
\label{sec:discussion}

\subsection{Five Promising Directions}

\begin{enumerate}
  \item \textbf{Feed-forward 4DGS.} Eliminate per-scene
        optimisation. L2D2-GS~\cite{song2026l2d2gs} opens this path;
        we expect end-to-end ``video in, 4DGS out'' models in the
        near future.
        \begin{itemize}
        \item \textit{Missing piece:} an open, multi-view dynamic
              dataset large enough to train a feed-forward 4DGS
              backbone comparable in scale to Objaverse-XL.
        \item \textit{What the next paper should look like:} a
              transformer-based densifier that consumes monocular
              video at 30~FPS and emits a 4DGS scene in one pass,
              evaluated on DyCheck~\cite{gao2022dycheck}.
        \item \textit{Mobile outlook:} feed-forward inference
              maps directly to the
              30~FPS sustained target expected on
              current Snapdragon-class mobile GPUs.
        \item \textit{Quantitative target:} by 2027~H1, a feed-forward
              4DGS backbone that reconstructs a 10-second DyCheck clip
              in $<$30~s on a single mobile GPU, matching the
              per-scene optimisation quality of Deformable-3DGS
              (LPIPS $\le 0.18$) within $5\%$.
        \end{itemize}

  \item \textbf{Hardware-software co-design.} Neo~\cite{oh2025neo}'s
        Reuse-and-Update sorting accelerator hints at dedicated
        4DGS compute units inside mobile GPUs.
        \begin{itemize}
        \item \textit{Missing piece:} a public 4DGS workload
              characterisation (FLOPS, tile occupancy, register
              pressure) on Adreno-class mobile GPUs.
        \item \textit{What the next paper should look like:} a
              tile-sort accelerator IP synthesised for a mobile
              shader core, benchmarked against the
              Vulkan~1.3 compute path.
        \item \textit{Mobile outlook:} the Adreno shader ISA and
              the Vulkan~1.3 compute-shader model are the
              candidate substrate for a co-designed sort
              accelerator.
        \item \textit{Quantitative target:} a co-designed sort
              accelerator that reduces the per-Gaussian sort cost
              by $3\times$ versus the software-only Vulkan~1.3
              baseline, measured on the Mip-NeRF~360 bicycle
              scene at 1080p.
        \end{itemize}

  \item \textbf{4DGS ray tracing.} The 4DGRT~\cite{liu20254dgrt}
        route is the long-term path to high fidelity, but mobile ray
        tracing remains an open problem.
        \begin{itemize}
        \item \textit{Missing piece:} mobile RT cores are
              bandwidth-starved for the per-Gaussian intersection
              workload; no published 4DGS BVH layout.
        \item \textit{What the next paper should look like:} a
              two-level BVH that exploits Gaussian covariance
              locality, evaluated against the rasteriser baseline
              on a Snapdragon-class mobile GPU.
        \item \textit{Mobile outlook:} mobile RT is not on the
              immediate Vulkan~1.3 hot path; keep as a 2027+
              watch item.
        \item \textit{Quantitative target:} by 2028, a BVH-based
              4DGS ray tracer that reaches 15~FPS @ 720p on
              a future Snapdragon (Gen~5 or later) mobile GPU,
              within $2\times$ of the rasteriser FPS. (This
              is a survey-side projection, not a
              4DGRT~\cite{liu20254dgrt} measurement; 4DGRT
              reports only desktop RTX~4090 / OptiX numbers,
              e.g.\ 36.56~FPS pinhole / 3.44~FPS fisheye, and
              does not project to mobile hardware.)
        \end{itemize}

  \item \textbf{Temporal compression.} Dynamic / static separation
        + residual coding + contextual coding may converge into a
        4DGS compression standard.
        \begin{itemize}
        \item \textit{Missing piece:} an MPEG-style codec profile
              that standardises the dynamic / static split, the
              residual format, and the rate-control interface.
        \item \textit{What the next paper should look like:} a
              reference encoder / decoder pair with bitstream
              conformance tests, integrating
              4DGS-CC~\cite{chen20254dgscc} and
              P-4DGS~\cite{wang2025p4dgs}.
        \item \textit{Mobile outlook:} a codec standard
              directly addresses the storage bottleneck
              observed in the surveyed mobile deployments.
        \item \textit{Quantitative target:} a reference codec that
              compresses the 4DGS-1K D-NeRF scenes to $<$20~MB
              per scene at PSNR $\ge 30$~dB, with a
              decoder-side decode time under 50~ms on
              a current mobile GPU.
        \end{itemize}

  \item \textbf{Generative 4DGS + cross-domain fusion.} The
        L2D2-GS~\cite{song2026l2d2gs} feed-forward route
        (\S\ref{sec:training}) generalises to a wider class of
        generative 4DGS models. The natural extension is
        text-to-4DGS (analogous to DreamFusion's text-to-3D path)
        and image-to-4DGS (single-image dynamic scene
        reconstruction). Cross-domain fusion closes the loop
        between 4DGS reconstruction and adjacent modalities
        (GaussianFluent~\cite{huang2026gaussianfluent} for
        physical simulation, 4DGS-SLAM for tracking, and
        on-device large multimodal models for grounding).
        \begin{itemize}
        \item \textit{Missing piece:} a public, large-scale
              text-paired dynamic video corpus (comparable in
              scale to Objaverse-XL or LAION-5B for static
              3D generation) and an end-to-end system paper
              that closes the loop between 4DGS reconstruction,
              physics simulation, and on-device LLM/VLM grounding.
        \item \textit{What the next paper should look like:} a
              diffusion-based 4DGS generator that consumes a
              text prompt + a single static 3DGS seed and
              outputs a 4DGS scene in $<$10~seconds on a
              current mobile GPU, evaluated on a held-out subset
              of DyCheck prompts; OR a ``video in,
              simulator-in-the-loop 4DGS out'' demo running in
              real time on a mobile device.
        \item \textit{Mobile outlook:} exciting but
              currently compute-bound; a mobile target
              is plausible only after the mobile
              transformer cost analysis matures.
        \item \textit{Quantitative target:} by 2027~H2, a
              text-to-4DGS system that produces a 10-second
              dynamic clip on a single mobile GPU in
              $<$30~seconds, with LPIPS $\le 0.20$ against the
              ground-truth clip; OR a cross-domain demo that
              combines 4DGS reconstruction with a
              fracture-aware simulation and renders the result
              at $\ge 10$~FPS on a mobile device.
        \end{itemize}

  \item \textbf{4DGS standardisation.} The 4DGS community
        currently has no standard codec profile, no standard
        benchmark suite, and no standard mobile evaluation
        harness. Three pieces of standardisation are likely
        in the next two years.
        \begin{itemize}
        \item \textit{Missing piece:} an MPEG- or
              Khronos-style codec profile that standardises the
              dynamic / static split, the residual format, and
              the rate--control interface. The
              4DGS-CC~\cite{chen20254dgscc} + P-4DGS~\cite{wang2025p4dgs}
              + PD-4DGS trio are the de-facto building
              blocks.
        \item \textit{What the next paper should look like:} a
              reference encoder / decoder pair with bitstream
              conformance tests, distributed as an open-source
              reference implementation.
        \item \textit{Mobile outlook:} the reference codec
              is the natural first deliverable for any
              mobile-4DGS working group.
        \item \textit{Quantitative target:} a reference
              decoder that runs at 30~FPS sustained on a
              Snapdragon-class mobile device, on a 20~MB 4DGS
              bitstream, with PSNR $\ge 30$~dB.
        \end{itemize}
\end{enumerate}

\subsection{Why Hardware--Software Co-Design Will Dominate}
We expect \emph{hardware--software co-design} to matter more than
algorithmic progress for mobile 4DGS over the next two years. The
argument has three legs.

\paragraph{First, the algorithmic headroom is shrinking.}
The 4DGS literature has seen roughly 3$\times$ per-year gains in
training speed (4DGS 30--60~min in 2023, OMG4 22~min in 2025,
L2D2-GS $<$1~s in 2026 at inference) and 5--10$\times$ per-year
gains in model size (vanilla 4DGS 60--200~MB in 2023,
4DGS-1K $\sim$7--50~MB (D-NeRF aggregate) / $\sim$50--420~MB
(N3V aggregate) in 2025, with per-scene ranges
D-NeRF 3.56--16.11~MB and N3V 29.34--69.24~MB).
Both curves are flattening: the next 3$\times$ in either dimension
will require a fundamentally new representation (e.g.\ generative
4DGS, see \S8.5 above), not an incremental tweak of the existing
one.

\paragraph{Second, the hardware headroom is large and quantifiable.}
The Adreno-class mobile GPU's tile-sort cost (10--20\% of the
frame budget at 60~FPS, see \S\ref{sec:rendering}) is a structural
artefact of the rasteriser's $\alpha$-blend rule
(Equation~\eqref{eq:alpha}). No algorithm change can amortise that
cost; only a hardware change can. Neo's Reuse-and-Update sort
(\S\ref{sec:rendering}) is a software prototype of what a
dedicated sort accelerator would do, and the 3--5$\times$ sort
speed-up it delivers would translate directly to a
30--50\% end-to-end speed-up at the 30~FPS sustained target.

\paragraph{Third, the co-design loop is short.}
The Snapdragon mobile launch cadence is roughly 12~months
(Gen~3 $\to$ Gen~4 $\to$ Gen~5 in 2024/2025/2026), which means
the algorithm-to-silicon loop is at most 24~months. The Adreno
shader ISA is stable across generations (no breaking changes
from Gen~3 to Gen~4), and the Vulkan~1.3 driver is upstreamed
into the Linux kernel with a 6-month cadence. These three
properties together make a tight algorithm / driver / ISA
co-design loop tractable, which is the structural advantage
that mobile 4DGS has over its desktop counterpart. The
contrarian view is therefore not ``algorithm doesn't matter''
but rather ``the next 3$\times$ on mobile will come from
hardware / driver / ISA co-design, not from a better MLP''.

\subsection{Broader Outlook}
The five directions above are not project-specific; they
generalise to any mobile-SoC target that combines a
tile-based deferred rasteriser with a Vulkan~1.3+ driver
(Qualcomm Snapdragon 7/8 series, MediaTek Dimensity 9000+,
Samsung Exynos 2400, etc.). The hardware-software co-design
thesis is a community-wide observation, not tied to any
single application. The Section~\ref{sec:mobile} quantitative
anchors (Flux-GS, Mobile-GS, GS-NFS, Lumina) serve as the
on-device measurement floor against which the directions
above should be evaluated.

\subsubsection{Additional related work surveyed for completeness.}
Beyond the 60 most-cited works discussed above, the present
survey also covers nine additional papers that inform the
future-directions discussion but did not fit naturally into
the in-line citations. They are cluster-cited here across
four thematic axes: \textbf{compression and dynamic/static
separation} (CATSplat~\cite{roh2025catsplat},
MaGS~\cite{ma2025mags}, LUDVIG~\cite{marrie2025ludvig},
SplatTalk~\cite{thai2025splattalk},
BezierGS~\cite{ma2025beziergs}); \textbf{self-calibrating
and self-supervised reconstruction}
(Self-Cali-GS~\cite{deng2025selfcalibratinggs},
GaussianOcc~\cite{gan2025gaussianocc}); \textbf{single-image
free-pose generation}
(FreeSplatter~\cite{xu2025freesplatter}). Each entry
serves as a borderline-or-out-of-faction reference that
expands the survey's coverage without diluting the four
core research directions.
